\section{Анализ предметной области}

\subsection{Особенности городских цифровых туристических сервисов}

Современный городской туристический сервис решает практическую задачу быстрого ориентирования пользователя в пространстве города: где находится интересный объект, как до него добраться и какие точки целесообразно объединить в одну прогулку. Для локального веб-проекта это особенно важно, поскольку аудитория обычно включает две группы с разными целями: жителей города и туристов.

Жители чаще ищут новые варианты досуга в знакомой среде: парки, кафе, культурные площадки, сезонные маршруты выходного дня. Туристы, напротив, ориентированы на компактный набор «ключевых» мест и ограничены временем пребывания. Следовательно, предметная область городского гида должна поддерживать одновременно точечный выбор объекта и сценарий последовательного посещения нескольких локаций.

Для проекта существенной особенностью является сочетание информационной и навигационной функций. Пользователь не только читает описание объекта, но и сопоставляет его с географическим положением на карте. Поэтому ценность сервиса определяется не объёмом текстов, а согласованностью данных: карточка места, категория, визуальные материалы и координаты должны описывать один и тот же объект без противоречий.

\subsection{Характеристика деятельности сервиса}

Деятельность городского гида строится как непрерывный цикл работы с контентом:
\begin{itemize}
	\item сбор и структурирование информации о местах города;
	\item публикация объектов в каталоге и на карте;
	\item объединение отдельных точек в тематические и прогулочные маршруты;
	\item регулярная актуализация сведений по мере изменений городской среды.
\end{itemize}

Таким образом, сервис выполняет роль локального агрегатора проверенной туристической информации. С практической точки зрения это означает, что предметная область включает не только сами объекты интереса, но и процессы их сопровождения: уточнение описаний, корректировку координат, обновление видимости карточек и поддержание маршрутов в актуальном состоянии.

В организационном плане выделяются два контура:
\begin{enumerate}
	\item Пользовательский --- поиск, просмотр, выбор мест и маршрутов.
	\item Административный --- ведение справочников и редактирование данных.
\end{enumerate}

Согласованность этих контуров напрямую влияет на качество сервиса: пользовательский интерфейс должен получать структурированные и актуальные данные, а контентный контур --- обеспечивать их своевременное обновление.

\subsection{Проблематика традиционного подхода к представлению туристической информации}

При отсутствии специализированной системы сведения о городских объектах обычно распределены по разным источникам: таблицам, текстовым документам, постам в социальных сетях и внешним картографическим сервисам. Такой подход приводит к ряду типовых проблем.

\textbf{Разрозненность данных.} Одни и те же объекты описываются в нескольких местах, что создаёт дублирование и расхождения в формулировках, адресах и характеристиках.

\textbf{Сложность актуализации.} Любое изменение (например, уточнение режима работы или статуса локации) требует ручной правки в нескольких каналах, из-за чего часть информации быстро устаревает.

\textbf{Высокие издержки пользовательского поиска.} Пользователь вынужден переключаться между источниками, чтобы собрать целостную картину: описание места --- в одном ресурсе, фото --- в другом, маршрут --- в третьем.

\textbf{Отсутствие связного маршрутного представления.} Даже при наличии перечня интересных точек не формируется готовый сценарий посещения, учитывающий порядок и логистику перемещения.

\textbf{Снижение доверия к сервису.} Неполные или несогласованные данные ухудшают пользовательский опыт и формируют ощущение ненадёжности ресурса.

Следовательно, традиционный «разрозненный» подход не соответствует потребностям городского туристического сервиса, где важны целостность информации, оперативность обновления и удобство навигации.

\subsection{Обоснование необходимости автоматизации процессов}

Автоматизация в рассматриваемой предметной области необходима не сама по себе, а как инструмент устранения перечисленных противоречий. Для городского гида она обеспечивает единое информационное пространство, в котором данные о месте, его категории, изображениях и координатах поддерживаются согласованно.

Практический эффект автоматизации проявляется в нескольких аспектах:
\begin{itemize}
	\item сокращается время обновления контента и публикации изменений;
	\item снижается вероятность ошибок при повторном ручном вводе;
	\item повышается доступность данных для конечного пользователя;
	\item упрощается сопровождение маршрутов как составных объектов;
	\item формируется единая логика представления городских локаций.
\end{itemize}

Таким образом, автоматизация является необходимым условием устойчивой работы сервиса, ориентированного на регулярное использование и постоянное развитие каталога мест.

\subsection{Структура предметной области и ключевые сущности}

На основании анализа деятельности сервиса выделяются следующие ключевые сущности:
\begin{enumerate}
	\item Категория --- классификация мест по типу (например, парк, музей, кафе, достопримечательность).
	\item Место --- конкретный объект городской среды с наименованием, описанием, адресом и координатами.
	\item Маршрут --- объединение нескольких мест в логическую последовательность посещения.
	\item Точка маршрута --- связь между маршрутом и местом, определяющая порядок прохождения.
\end{enumerate}

Данная структура отражает естественную логику пользовательского поведения: от общего интереса к типу досуга --- к выбору конкретных мест --- к формированию маршрута. Важным является то, что маршруты в этой модели выступают не вторичным списком, а самостоятельным объектом предметной области, имеющим собственную ценность для пользователя.

\subsection{Анализ существующих решений в тематике городских гидов}

Цифровая трансформация туристической отрасли — один из главных драйверов её развития в XXI веке. Сегодня недостаточно просто иметь перечень достопримечательностей на сайте: путешественник ожидает бесшовного пути от вдохновения к покупке билета, бронированию гостиницы и построению маршрута прямо в смартфоне. Именно поэтому региональные туристические платформы всё чаще превращаются в полноценные цифровые экосистемы, объединяющие информацию, сервисы, местный бизнес и самих жителей.
\subsubsection{Швейцарский подход}

Наиболее ярким примером инновационного подхода на национальном уровне является запуск Швецией первой в мире открытой туристической API (Application Programming Interface). Её опыт отличается системным подходом: это не единый государственный сервис, а экосистема кооперирующих платформ, каждая из которых отвечает за свой сегмент данных. Рассмотрим архитектуру этой системы подробнее.
В отличие от модели единого национального портала, швейцарский подход строится на взаимодействии нескольких ключевых элементов, объединённых общими стандартами.

\textbf{Discover.swiss} --- технологическое ядро отрасли. Это не просто API, а полноценный Knowledge Graph («цифровой двойник швейцарского туризма»), который связывает информацию о местах, событиях, турах, погоде и веб-камерах в единую семантическую сеть. Ключевая особенность — данные хранятся в формате schema.org и JSON-LD, что обеспечивает совместимость с поисковыми системами и сторонними приложениями. К системе уже подключены десятки внешних источников: системы бронирования столиков (Aleno, Lunchgate), платформы активного туризма (Outdooractive), системы управления горнолыжными курортами, данные по доступной среде (Sitios) и PMS-системы отелей.

\textbf{Accommodation Data Hub} --- отраслевое решение для средств размещения. Проект, профинансированный SECO (Государственный секретариат по экономике) с грантом 281 689 CHF, решает проблему фрагментации данных о гостиничной базе. Отельеры через единый интерфейс myAccommoData управляют маркетинговой информацией (фото, описания удобств, услуги), которая затем автоматически обогащается данными о классификации («звёзды»), наличии экомаркировок и доступности; бесплатно распространяется через API discover.swiss для всех участников рынка — от национального офиса Switzerland Tourism до региональных DMO.

\textbf{Open Journey Planner (OJP)} --- мультимодальная мобильность.
API планировщика маршрутов, работающий «от координаты до координаты», позволяет строить поездки с использованием всех видов транспорт. Сервис обслуживается платформой opentransportdata.swiss и доступен по API-ключу. OJP критически важен для туристических приложений, поскольку связывает объекты размещения и достопримечательности с транспортной сетью страны.

\textbf{Метакаталог Tourism Data Catalogue}. Портал выполняет роль «единого окна» для разработчиков, агрегируя метаданные обо всех доступных источниках туристической информации в стране. Здесь можно найти документацию по API достопримечательностей, контактные данные провайдеров и условия использования.

Если в России региональные платформы часто идут по пути создания замкнутых экосистем «под ключ» (единый агрегатор с функцией бронирования и эквайрингом), то швейцарская модель ориентирована на инфраструктуру обмена данными. Государство или крупная ассоциация создаёт API и стандарты, а частный бизнес конкурирует в создании приложений на этой основе. Это требует высокого уровня зрелости рынка и доверия между игроками, но обеспечивает более быстрое масштабирование инноваций.

Общеевропейский тренд на кооперацию подтверждается инициативами в рамках программ Interreg Central Europe. В проекте Capacity2Transform участвуют девять регионов, создающих цифровые решения для устойчивого туризма: от платформы для велотуризма в Словении до мультисенсорных туристско-информационных центров в Чехии. Это свидетельствует о понимании на уровне ЕС того факта, что цифровая платформа сегодня — не маркетинговый буклет, а инструмент управления региональным развитием.
\subsubsection{Региональные платформы Российской Федерации}

В Российской Федерации развитие региональных платформ идёт по пути создания комплексных экосистем, решающих задачи не только информирования туриста, но и прямого бронирования услуг, а также глубокой аналитики турпотока.

\textbf{«Курортный помощник Гоша»} --- проект, стартовавший как локальное решение для города Сочи, к 2025 году масштабировался до федеральной сети, охватывающей восемь туристических регионов России. В отличие от государственных информационных порталов, данная платформа представляет собой частную экосистему маркетплейс-типа, в основе которой лежит экономическая модель win-win для трёх ключевых групп стейкхолдеров: туристов, поставщиков услуг и отелей.
Сервис разрабатывался как ответ на комплекс системных проблем, характерных для фрагментированного рынка туристских услуг российских курортов. Исследование болей целевой аудитории, проведённое разработчиками, выявило следующие барьеры:
\begin{enumerate}
	\item Для туристов: трудоёмкость поиска и сравнения досуговых мероприятий из-за разрозненности информации на десятках сайтов; непрозрачное ценообразование и риск переплаты перекупщикам; неудобство хранения билетов из разных каналов продаж; сложность оформления возвратов.
	\item Для поставщиков услуг (экскурсионные бюро, организаторы активностей): отсутствие единой цифровой площадки для продвижения; ведение бронирований вручную (Excel, мессенджеры, бумажные журналы); высокие маркетинговые затраты без гарантии конверсии; риск незаполненности мест в высокий сезон.
	\item Для отелей и гостевых домов: потеря потенциального дохода от рекомендаций гостям (туристы покупают услуги на стороне) и отсутствие удобного инструмента для монетизации собственной экспертизы о досуге в регионе.
\end{enumerate}

Платформа построена по модульному принципу, где каждый тип пользователя получает специализированный личный кабинет с набором релевантных инструментов.

Механика взаимодействия: ключевым звеном экосистемы выступает отель. При заселении турист получает QR-код, ведущий на регистрацию в сервисе, что автоматически привязывает гостя к конкретному средству размещения. В дальнейшем любая покупка, совершённая туристом через платформу, генерирует комиссионный доход, распределяемый между сервисом, поставщиком услуги и рекомендовавшим отелем.

\textbf{Портал «КАК ЕСТЬ»} --- это уникальный для России пример туристической экосистемы, основанной на принципе community-driven, то есть на контенте, который создают сами пользователи. В отличие от многих официальных гидов-справочников, этот проект изменил подход к внутреннему туризму в Самарской области, сделав ставку на живые впечатления, интерактивность и локальную аутентичность. На платформе есть такие разделы, как афиша мероприятий, фестивалей, выставок и концертов; обширный каталог локаций — от культовых достопримечательностей до секретных природных точек; тотовые сценарии для прогулок и путешествий на любой вкус и время; а также представлены эко-отели - это самая полная база из более чем 150 загородных объектов для размещения в регионе - от кемпингов до глэмпингов. На платформе есть субъективный, но очень влиятельный рейтинг ресторанов, глэмпингов и других мест. Редакционная команда лично посещает заведения, чтобы оценить кухню, сервис и атмосферу. Это создает высокий уровень доверия у пользователей и даже стало маркетинговым инструментом для местного бизнеса, так как позиции в рейтинге могут меняться после каждого нового визита экспертов. 

\textbf{Russia.Travel} --- это не просто сайт с описаниями достопримечательностей, а ключевой элемент государственной стратегии по цифровизации туризма. С конца 2023 года это единая экосистема, которая объединяет вдохновляющий контент о путешествиях по России с практическими инструментами для их планирования, став главным цифровым окном в мир внутреннего туризма. Свой путь Russia.Travel начал в 2013 году. Однако ключевой момент в его истории — масштабный перезапуск 28 декабря 2023 года. Russia.Travel объединили с московской цифровой платформой RUSSPASS. Это слияние легло в основу единой цифровой туристической экосистемы России. Данная платформа как «витрина впечатлений»: ее задача — вдохновлять на путешествия с помощью статей, фотографий и описаний, представляя весь туристический потенциал страны. RUSSPASS как «технологическое ядро», он предоставляет практические инструменты для планирования, являясь полноценным туристическим маркетплейсом. На сайте доступна обширная база, включающая более 50 национальных и свыше 150 маршрутов выходного дня. Они удобно распределены по 7 географическим зонам, для каждого маршрута есть подробный путеводитель. Можно покупать билеты на мероприятия (включая оплату «Пушкинской картой»), бронировать отели, экскурсии и другие услуги через экосистему RUSSPASS. Также включает фотобанк с более чем 3 000 бесплатных изображений из 33 регионов, а также статьи и анонсы событий. Анализ поведения пользователей на платформе — ключ к ее непрерывному улучшению и принятию верных стратегических решений по развитию туризма в регионе.
 
\textbf{Культурно-туристический портал «Соловьиный край»} --- это главная онлайн-площадка (проект регионального Комитета по культуре и компании «Ростелеком») обо всех культурных и туристических событиях Курской области. Портал был запущен 16 декабря 2021 года, чтобы объединить на одном ресурсе максимум информации о мероприятиях и достопримечательностях региона. На сайте портала можно посмотреть афишу мероприятий, а также приобрести билеты (здесь собрана полная афиша концертов, спектаклей и других культурных событий, на которые можно купить билеты онлайн без наценки). На сайте представлены туристические маршруты - для путешественников разработаны готовые маршруты по главным достопримечательностям, таким как Коренная пустынь, усадьба Афанасия Фета, Парк мельниц, а также, к примеру, маршрут «Соловьи и железо». Также можно подобрать гостиницу, найти интересные кафе. Также доступна опция подписки на новостную рассылку с персональными рекомендациями. Бонусом является то, что портал активно участвует в программе, позволяя молодым людям в возрасте от 14 до 22 лет посещать события за счет государства.

\textbf{«Go Kursk»} --- это не один сервис, а целая экосистема туристических продуктов Курской области. Её основа -- Туристско-информационный центр (ТИЦ) Курской области, который объединяет сайт, соцсети и мобильные приложения, помогая гостям и жителям региона. Основной сайт -- это главный информационный ресурс, своего рода «точка входа» для планирования путешествия. На портале, доступном в том числе на английском и немецком языках, собрана информация об истории, достопримечательностях, гостиницах, кафе и событиях. Сайт предлагает готовые маршруты, например, национальный «Соловьи и железо» или «Курский соловей» на поезде, а также помогает подобрать экскурсии и забронировать проживание. Кроме того, он активно участвует в программе «Пушкинская карта». Для удобства навигации прямо во время путешествия созданы специальные приложения, которые можно найти в RuStore: «Курский путь»: приложение-гид по досугу непосредственно в городе Курске. Помогает найти интересные события, вы ставки и фестивали по вашим интересам, а также содержит подборку любопытных фактов о городе. Существует даже краеведческая игра, созданная для привлечения туристов. Она позволяет в интерактивной форме изучить историю края.

\textbf{Платформа «Я — курянин»} --- это уникальный культурно-просветительский проект с сильным туристическим уклоном. По своей направленности он отличается от сервисов, которые были описаны ранее, предлагая не просто афишу или готовый маршрут, а глубокое погружение в историю и культуру региона через исследование и обучение. Это не просто туристический сервис, а часть региональной Стратегии развития образования. Платформа служит виртуальной средой, созданной педагогами и учёными для детей, родителей и учителе. Платформа представляет собой интерактивную карту-конструктор. Вы отмечаете интересные вам категории, а система сама подбирает релевантную информацию: от культурных событий и архитектурных памятников до сведений о выдающихся земляках. Контент платформы очень разнообразен и упорядочен по тематическим маршрутам: многообразие категорий, среди которых можно найти «Туристические маршруты», «История курского края», «Выдающиеся земляки», а также «Культурные традиции», «Гастрономические ремесленные маршруты» и «Духовно-нравственное воспитание» (включая монастыри и храмы). пользователи также могут изучать «Археологические», «Литературные», «Искусствоведческие» и «Военно-патриотические» направления. На сайте также имеются более 400 цифровых открыток. Это визуальные карточки, созданные учениками курских школ, посвящены самым разным достопримечательностям региона и наполняют «Атлас путешествий».
\subsubsection{Выводы по анализу существующих решений}

Проведенный анализ позволяет выделить несколько ключевых принципов, на которых строятся успешные современные туристические платформы. Самые успешные платформы стремятся объединить максимум услуг в одном окне, решая проблему фрагментации информации для пользователя. Веб-приложение — это уже не опция, а базовый элемент, обеспечивающий доступ к информации и сервисам в любой момент путешествия. Пользователям нужны не просто списки достопримечательностей, а curated experiences — тщательно отобранные впечатления и маршруты от местных экспертов. Важно создавать сервисы, удобные для всех, включая людей с ограниченными возможностями, и снимающие барьеры для иностранных туристов.



